% =========================================================
% Real-Valued Functions — Notes
% Chapter: functions
% Section: real-valued-functions
% Volume III · Analysis
% =========================================================

\section{Real-Valued Functions}
\label{sec:real-valued-functions}

% ---------------------------------------------------------
% Toolkit
% ---------------------------------------------------------

\begin{tcolorbox}[title={Toolkit --- Real-Valued Functions}]
\begin{itemize}
  \item Pointwise operations on $f, g : X \to \mathbb{R}$: sum,
    difference, product, quotient (where $g \neq 0$), scalar
    multiple, absolute value, max, min. Each is purely algebraic —
    no limits involved.
  \item $fg$ denotes pointwise product $(fg)(x) := f(x)g(x)$, not
    composition $f \circ g$.
  \item Bounded above: $\exists M, \forall x, f(x) \leq M$.
    Bounded: $\exists B > 0, \forall x, |f(x)| \leq B$.
  \item Increasing: $x_1 < x_2 \Rightarrow f(x_1) \leq f(x_2)$.
    Strictly increasing: $x_1 < x_2 \Rightarrow f(x_1) < f(x_2)$.
  \item Maximum attained at $x^*$: $f(x^*) = \sup f(A)$ with $x^* \in A$.
    Supremum need not be attained.
\end{itemize}
\end{tcolorbox}

% ---------------------------------------------------------
% Definition — Pointwise Arithmetic Operations
% ---------------------------------------------------------

\begin{tcolorbox}[title={Definition --- Pointwise Arithmetic Operations}]
\begin{definition}[Pointwise Arithmetic Operations on Real-Valued Functions]
\label{def:pointwise-ops-R}
Let $X \subseteq \mathbb{R}$ and let $f, g : X \to \mathbb{R}$.
The following real-valued functions on $X$ are defined pointwise:
\begin{align*}
  (f + g)(x)      &:= f(x) + g(x)                         & &\text{sum} \\
  (f - g)(x)      &:= f(x) - g(x)                         & &\text{difference} \\
  (fg)(x)         &:= f(x) \cdot g(x)                     & &\text{product} \\
  (cf)(x)         &:= c \cdot f(x),\quad c \in \mathbb{R} & &\text{scalar multiple} \\
  \max(f,g)(x)    &:= \max(f(x), g(x))                    & &\text{pointwise maximum} \\
  \min(f,g)(x)    &:= \min(f(x), g(x))                    & &\text{pointwise minimum} \\
  |f|(x)          &:= |f(x)|                               & &\text{absolute value}
\end{align*}
If additionally $g(x) \neq 0$ for all $x \in X$:
\[
  (f/g)(x) := f(x)/g(x). \qquad\text{quotient}
\]
\end{definition}
\end{tcolorbox}

\begin{remark*}[Standard quantified statement]
\[
  \forall x \in X :\; (f+g)(x) = f(x) + g(x).
\]
Each operation is pointwise: the value at $x$ depends only on $f(x)$
and $g(x)$, not on values at nearby points.
\end{remark*}

\begin{remark*}[Product vs.\ composition]
$fg : X \to \mathbb{R}$ maps $x \mapsto f(x)g(x)$. This is the
pointwise product, not the composition $f \circ g$ which maps $x
\mapsto f(g(x))$ and requires $g : X \to X$. For $f(x) = x^2$ and
$g(x) = 2x$: $(fg)(x) = 2x^3$ but $(f \circ g)(x) = 4x^2$.
\end{remark*}

\begin{remark*}[Algebraic structure]
The set $\mathbb{R}^X := \{f : X \to \mathbb{R}\}$ forms a
commutative ring under pointwise addition and multiplication, with
zero function $\mathbf{0}(x) := 0$ and identity $\mathbf{1}(x) :=
1$. Scalar multiplication makes it an $\mathbb{R}$-algebra.
\end{remark*}

% ---------------------------------------------------------
% Definition — Bounded Functions
% ---------------------------------------------------------

\begin{tcolorbox}[title={Definition --- Boundedness of Functions}]
\begin{definition}[Bounded Above, Bounded Below, Bounded]
\label{def:function-boundedness}
Let $A \subseteq \mathbb{R}$ and let $f : A \to \mathbb{R}$.
\begin{itemize}
  \item $f$ is \textbf{bounded above} on $A$ if $\exists M \in
    \mathbb{R},\; \forall x \in A :\; f(x) \leq M$.
  \item $f$ is \textbf{bounded below} on $A$ if $\exists m \in
    \mathbb{R},\; \forall x \in A :\; f(x) \geq m$.
  \item $f$ is \textbf{bounded} on $A$ if $\exists B > 0,\;
    \forall x \in A :\; |f(x)| \leq B$.
\end{itemize}
\end{definition}
\end{tcolorbox}

\begin{remark*}[Standard quantified statement]
\[
  f \text{ bounded on } A
  \;\iff\;
  \exists B > 0,\;
  \forall x \in A :\; |f(x)| \leq B
  \;\iff\;
  f(A) \subseteq [-B, B].
\]
\end{remark*}

\begin{remark*}[Negated quantified statement]
\[
  f \text{ unbounded on } A
  \;\iff\;
  \forall B > 0,\;
  \exists x \in A :\; |f(x)| > B.
\]
\end{remark*}

\begin{remark*}[Interpretation]
Bounded means the image $f(A)$ fits inside some finite interval
$[-B, B]$. Bounded above alone does not imply bounded: $f(x) := -x$
on $\mathbb{R}$ is bounded above by $0$ but unbounded below.
\end{remark*}

% ---------------------------------------------------------
% Definition — Supremum, Infimum, Maximum, Minimum
% ---------------------------------------------------------

\begin{tcolorbox}[title={Definition --- Extremal Values of Functions}]
\begin{definition}[Supremum, Infimum, Maximum, Minimum of a Function]
\label{def:function-sup-inf-max-min}
Let $A \subseteq \mathbb{R}$ and let $f : A \to \mathbb{R}$.
\begin{itemize}
  \item The \textbf{supremum} of $f$ on $A$ is $\sup_{x \in A}
    f(x) := \sup f(A)$.
  \item The \textbf{infimum} of $f$ on $A$ is $\inf_{x \in A}
    f(x) := \inf f(A)$.
  \item $f$ \textbf{attains a maximum} at $x^* \in A$ if
    $f(x^*) = \sup_{x \in A} f(x)$.
  \item $f$ \textbf{attains a minimum} at $x_* \in A$ if
    $f(x_*) = \inf_{x \in A} f(x)$.
\end{itemize}
\end{definition}
\end{tcolorbox}

\begin{remark*}[Standard quantified statement]
\[
\begin{aligned}
  f \text{ has max at } x^*
  &\;\iff\;
  x^* \in A \;\wedge\; \forall x \in A :\; f(x) \leq f(x^*). \\
  f \text{ has min at } x_*
  &\;\iff\;
  x_* \in A \;\wedge\; \forall x \in A :\; f(x) \geq f(x_*).
\end{aligned}
\]
\end{remark*}

\begin{remark*}[Interpretation]
The supremum of $f$ on $A$ always exists when $f$ is bounded above
(by completeness), but need not be attained. A maximum is a
supremum that is achieved at some point of the domain. The function
$f(x) = x$ on $(0,1)$ has supremum $1$ but no maximum: no $x \in
(0,1)$ achieves $f(x) = 1$.
\end{remark*}

% ---------------------------------------------------------
% Definition — Monotone Functions
% ---------------------------------------------------------

\begin{tcolorbox}[title={Definition --- Monotonicity}]
\begin{definition}[Increasing, Decreasing, Monotone, Constant]
\label{def:function-monotonicity}
Let $f : A \to \mathbb{R}$ on a nonempty $A \subseteq \mathbb{R}$.
\begin{itemize}
  \item $f$ is \textbf{increasing} on $A$ if $x_1 < x_2 \Rightarrow
    f(x_1) \leq f(x_2)$.
  \item $f$ is \textbf{strictly increasing} on $A$ if $x_1 < x_2
    \Rightarrow f(x_1) < f(x_2)$.
  \item $f$ is \textbf{decreasing} on $A$ if $x_1 < x_2 \Rightarrow
    f(x_1) \geq f(x_2)$.
  \item $f$ is \textbf{strictly decreasing} on $A$ if $x_1 < x_2
    \Rightarrow f(x_1) > f(x_2)$.
  \item $f$ is \textbf{monotone} if it is increasing or decreasing.
  \item $f$ is \textbf{constant} if $f(x_1) = f(x_2)$ for all
    $x_1, x_2 \in A$.
\end{itemize}
\end{definition}
\end{tcolorbox}

\begin{remark*}[Standard quantified statement]
\[
\begin{aligned}
  f \text{ increasing}
  &:\; \forall x_1, x_2 \in A,\; x_1 < x_2 \Rightarrow f(x_1) \leq f(x_2). \\
  f \text{ strictly increasing}
  &:\; \forall x_1, x_2 \in A,\; x_1 < x_2 \Rightarrow f(x_1) < f(x_2).
\end{aligned}
\]
\end{remark*}

\begin{remark*}[Interpretation]
The word \emph{strictly} tightens $\leq$ to $<$. A constant function
is both increasing and decreasing but neither strictly increasing nor
strictly decreasing. A strictly increasing function is injective:
$f(x_1) = f(x_2) \Rightarrow x_1 = x_2$ (by contrapositive of the
strict condition).

Terminology is not uniform across sources. These notes follow Bartle
and Sohrab: \emph{increasing} means non-decreasing; \emph{strictly
increasing} means the strict inequality holds.
\end{remark*}
